% Korean Noether paper 37 — fresh producer translation, unchecked.
% Initial translation authority: ED0005 lines 18639--18830.
% Current scope-rebase authority: ED0011 / NOETH-DE-AUTH-v047-20260808.
% Authority SHA-256: 1A44F967B29972E8F99E5C323A479162AD82A23FC457395915A4BB9DDF51AD41.
% Normalized source span: 22,155 bytes; SHA-256 11019FA922A13EF8915F002352E5ACFC3A8CD36AA3F65A856DC12B8A7F4CE973.
% Translation only; independent source, Korean-language, render, and certification checks remain external.

\section*{37. 고차 분기가 없는 체의 정규기저}

\begin{center}
\emph{Journal f. d. reine u. angew. Math. 167 (1932), 147--152쪽}
\end{center}

\begin{center}
괴팅겐의 \emph{에미 뇌터}.
\end{center}

갈루아 수체 \(K/k\)의 정규기저란 \(k\)의 정수환 \(\frako\)에 대한 \(K\)의 정수환 \(\frakO\)의 기저로서 켤레원소들로 이루어진 것을 말한다. 그러한 기저가 존재하기 위한 필요조건은 \(K/k\)에 고차 분기체가 나타나지 않는다는 것임이 알려져 있다.\footnote{A. Speiser, \emph{Gruppendeterminante und Körperdiskriminante}, § 6, Math. Ann. 77 (1916), 546--562쪽을 보라.} 이는 한 소에 국한하여, 곧 \(k\)의 \(\frakp\)-진 확대 \(k_\frakp\)와 이에 대응하는 \(K\)의 확대 \(K_\frakp\)로 넘어가도 성립한다. 이하에서는 그 역이 성립함을 보인다.

\emph{\(\frakp\)가 \(K/k\)의 차수를 나누지 않는 모든 소 \(\frakp\)에서 정규기저가 존재한다. 특히 \(K/k\)에 고차 분기가 없으면 모든 분기소에서 정규기저가 존재하며, 그곳에서 판별식은 군행렬식의 제곱이 된다.}

이 마지막 명제에는 다음 설명을 덧붙여야 한다. E. Artin의 추측과 달리 그의 도체들\footnote{E. Artin, \emph{Über die gruppentheoretische Struktur der Diskriminante}, J. f. Math. 164 (1931), 1--11쪽.}로 넘어가려면 더 많은 고찰이 필요하다. 군행렬식을 기약지표에 따라 분해하면 일반화된 근수들이 생기고, 이를 알맞게 묶으면 지표들을 첨가한 기초체에서의 분해를 얻는다. 가장 간단한 경우들에는 근수들의 아이디얼론적 분해를 이용하여 이 새 인자들을 Artin 도체와 동일시할 수 있다.\footnote{[1931. 10. 6.] 일반적인 경우에도 아이디얼론적 고찰이 더 필요함은 M. Deuring이 알려 준 다음 예에서 드러난다. 이 예는 임의로 변형할 수 있다. 어떤 비극대 오더의 판별식은 군행렬식의 제곱으로 나타나지만 켤레 지표들에는 서로 다른 아이디얼인자들이 대응한다. \(k\)를 5차 단위근체, \(K=k(\sqrt[5]{2})\)라 하자. 오더 \([1,2\sqrt[5]{2},\sqrt[5]{2^2},\sqrt[5]{2^3},\sqrt[5]{2^4}]\)를 생각하되, 이 글의 고찰에 따라 이 오더가 정규기저를 갖는 소 \((2)\)에서 생각한다. 지표들에 따른 대응 군행렬식의 분해에서는 위의 기저원소들이 바로 인자들로 나온다. 따라서 \glqq{}도체\grqq{}들은 \(1\)을 제외하면
\[
2\sqrt[5]{2}\sqrt[5]{2^4},\quad
\sqrt[5]{2^2}\sqrt[5]{2^3},\quad
\sqrt[5]{2^3}\sqrt[5]{2^2},\quad
\sqrt[5]{2^4}\,2\sqrt[5]{2},
\]
곧 \(4,2,2,4\)가 되어 켤레 지표들에 대해 서로 다르다.}
더 나아가 정규기저가 존재하지 않는 일반적인 경우에도 \glqq{}갈루아 가군에 따른 합성열\grqq{}을 이용하면 일반화된 근수들로 언제나 분해할 수 있음을 말해 두고자 한다. 여기에서도 위와 마찬가지로 확대된 기초체에서의 분해가 생기며, 다시 아이디얼론적 문제가 시작된다.

위 가정 아래에서 정규기저의 존재를 증명하는 근거는 다음과 같다. \(\frakO/\frako\)를 갈루아 가군, 곧 \(\frako\)-가군으로 보되 갈루아군의 치환들을 연산자영역으로 두면, 이는 \(\frako\)로 확대된 정수군환의 한쪽 아이디얼과 연산자 동형사상이 된다. 이 정수군환은 모든 보통 분기소에서 극대 오더이다. 그런데 \(\frakp\)-진 반단순계의 극대 오더에 속하는 아이디얼들은 주아이디얼이다.\footnote{H. Hasse, \emph{Über \(\frakp\)-adische Schiefkörper und ihre Bedeutung für die Arithmetik hyperkomplexer Zahlsysteme}, 정리 46, 47, Math. Ann. 104 (1931), 495--534쪽을 보라. 가환계에서는 \(k\)의 \(\frakp\)에 대한 국소화환으로 넘어가기만 해도 주아이디얼 성질이 성립한다. 그러므로 아벨군과 아벨체를 다룰 때에는 이 더 약한 확대를 써서 소를 정의해도 된다.} 따라서 여기에서는 한 원소에 군 또는 군환의 치환들을 적용하여 아이디얼을 얻는다. 갈루아 가군으로 되돌리면 이것이 바로 정규기저의 존재를 뜻한다. 고차 분기소에서는 정수군환이 비극대 오더가 되고 \(\frakO/\frako\)에 대응하는 아이디얼은 비주아이디얼이 된다.

상수체의 표수가 \(K/k\)의 차수를 나누지 않는 한, 이 모든 고찰은 한 변수의 함수체에도 그대로 성립한다.

\subsection*{§1. \(\frakp\)-진으로 확대한 정수군환}

\(\frako\)와 \(\frako_\frakp\)는 각각 대수적 수체 \(k\)와 그 \(\frakp\)-진 확대 \(k_\frakp\)의 정수환이며, \(\frakp\)는 \(k\)의 소아이디얼이라고 하자. 원소 수가 \(n\)인 군 \(\Gg\)의 유리수계수 군환을 \((\Gg)\), 정수군환을 \([\Gg]\)로 쓴다. \((\Gg)_k\)와 \((\Gg)_{k_\frakp}\)는 각각 \(k\)와 \(k_\frakp\)에 계수를 둔 군환을 뜻하고, \([\Gg]_\frako\)와 \([\Gg]_{\frako_\frakp}\)는 정수군환을 각각 \(\frako\)와 \(\frako_\frakp\)에 계수를 둔 군환으로 확대한 것을 뜻한다.

따라서 \((\Gg)_k\)와 \((\Gg)_{k_\frakp}\)는 라디칼이 없는 계, 곧 반단순계이고, \([\Gg]_\frako\)와 \([\Gg]_{\frako_\frakp}\)는 각각 그 안의 오더이다.

\paragraph{정리 1.}
\emph{소아이디얼 \(\frakp\)가 \(\Gg\)의 원소 수 \(n\)을 나누지 않으면 \([\Gg]_{\frako_\frakp}\)는 \((\Gg)_{k_\frakp}\)의 극대 오더이다. \(\frakp\)가 \(n\)을 나누면 \([\Gg]_{\frako_\frakp}\)는 극대가 아니다.}

정수군환의 판별식, 따라서 \([\Gg]_\frako\)와 \([\Gg]_{\frako_\frakp}\)의 판별식은 원소 수 \(n\)의 한 거듭제곱과 같다.\footnote{예를 들어 E. Noether, \emph{Hyperkomplexe Größen und Darstellungstheorie}, § 26, Math. Ztschr. 30 (1929), 641--692쪽을 보라.} 그러므로 \(\frakp\)가 \(n\)을 나누지 않으면 \([\Gg]_{\frako_\frakp}\)의 판별식은 단위원이다. 이는 \(\frako_\frakp\)를 고정한 채 \([\Gg]_{\frako_\frakp}\)를 더 큰 오더로 확대할 수 없다는 뜻이다. 그러한 확대가 존재하면 판별식에서 비단위원인 제곱인자가 떨어져 나오기 때문이다. 그런데 \(\frako_\frakp\)는 이미 \(k_\frakp\)의 정수환, 따라서 그 안의 극대 오더라고 가정하였다. 이로써 이하에서 유일하게 쓰일 정리 1의 첫 부분이 증명되었다.

반대로 \(\frakp\)가 \(n\)을 나누면 자명표현에 대응하는 직합
\[
\mathfrak C=
E^{(1)}[\Gg]_{\frako_\frakp}
+(1-E^{(1)})[\Gg]_{\frako_\frakp},
\qquad
E^{(1)}=\frac1n\sum S
\]
은 \([\Gg]_{\frako_\frakp}\)의 진확대이다. 실제로 \(E^{(1)}=\frac1n\sum S\)이므로 \(\mathfrak C\)는 진확대이며, \(E^{(1)}\)과 \(S\in\Gg\)에 대한 \((1-E^{(1)})S\)를 종속인 생성원들로 갖는 오더이다. 여기서 \(E^{(1)}S=E^{(1)}\)이다.

\paragraph{정리 2.}
\emph{\(\frakp\)가 \(n\)을 나누지 않으면 \([\Gg]_{\frako_\frakp}\)에 관한 모든 정수 또는 분수 아이디얼은 주아이디얼이다.}

정리 1에 따라 \([\Gg]_{\frako_\frakp}\)는 \(\frakp\)-진 반단순계의 극대 오더이므로 그 아이디얼들은 주아이디얼이다.\footnote{Hasse, 앞의 논문. 그곳에서는 단순계에 대해서만 주아이디얼 성질을 증명하지만 알려진 추론에 따라 반단순계에도 성립한다. 극대 오더의 성분들은 다시 극대이고, 따라서 생각하는 아이디얼의 성분들은 이를테면 기저 \(a_i\)를 가진 주아이디얼이다. 그러면 \(a=\sum a_i\)가 본래 아이디얼의 기저가 된다. 아벨군에 관해서는 주 3)도 보라.}

\subsection*{§2. 갈루아 가군, 연산자 동형사상, 정규기저}

\(K/k\)가 갈루아이고 그 군이 \(\Gg\)라고 하자. \(z^S\)는 \(S\in\Gg\)라는 치환으로 \(z\)에서 얻는 \(K\)의 원소를 뜻한다. \(\Gg\)의 치환들은 \(K/k\)를 \(k\)-가군, 곧 \(k\)를 연산자영역으로 갖는 가법 아벨군으로 볼 때 연산자 자기동형사상들을 만든다. 따라서 다음과 같이 정하면 \(K/k\)를 \((\Gg)_k\)에 관한 가군으로 볼 수 있다.
\[
z\Bigl(\sum_i S_i c_i\Bigr)=\sum_i z^{S_i}c_i,
\qquad z\in K,\quad S_i\in\Gg,\quad c_i\in k.
\]

\paragraph{정의.}
\emph{\(K/k\) 안의 모든 \((\Gg)_k\)-가군을 유리 갈루아 가군이라 한다. 마찬가지로 \(K/k\) 안의 \([\Gg]_\frako\)-가군들을 정수 갈루아 가군이라 한다. 특히 \(\frakO/\frako\)는 정수 갈루아 가군이다.}

\paragraph{정리 3.}
\emph{유리 갈루아 가군으로서 \(K/k\)는 \((\Gg)_k\)와 연산자 동형사상이다.}

이는 \(K/k\)에 언제나 체의 정규기저, 곧 켤레원소 \(z^S\)들로 이루어진 \(k\)-기저가 존재한다는 사실에서 따라온다.\footnote{\(a_1,\ldots,a_n\)이 \(K/k\)의 한 기저이고 \(u_i\)가 부정원들이면, \(S\)가 \(\Gg\)의 원소들을 돌 때 \(\sum u_i a_i^S\)들은 선형독립이다. 따라서 선형독립성을 보존하도록 \(u_i\)들을 \(k\)의 원소들로 특수화할 수도 있다.} 대응
\[
S\longmapsto z^S,\qquad
\sum_i S_i c_i\longmapsto\sum_i z^{S_i}c_i
\quad(c_i\in k),\qquad
ST\longmapsto z^{ST}
\]
은 연산자 준동형사상을 주고, 양변의 \(k\)-랭크가 같으므로 동형사상이 된다.

\paragraph{정리 4.}
\emph{정수 갈루아 가군으로서 \(\frakO/\frako\)는 \((\Gg)_k\) 안의 최고랭크 \(n\)인 한 \([\Gg]_\frako\)-가군과 연산자 동형사상이다. 마찬가지로 \(\frakO_\frakp/\frako_\frakp\)는 \((\Gg)_{k_\frakp}\) 안의 랭크 \(n\)인 한 \([\Gg]_{\frako_\frakp}\)-가군과 연산자 동형사상이다.}

정리 3의 \((\Gg)_k\)-동형사상 \(K/k\to(\Gg)_k\)는 \([\Gg]_\frako\)-가군 \(\frakO/\frako\)에 \((\Gg)_k\) 안의 랭크 \(n\)인 한 \([\Gg]_\frako\)-가군을 대응시킨다. 또 정리 3의 대응에서 \(c_i\)들이 \(k_\frakp\)의 모든 원소를 돌게 하면, \((\Gg)_k\)-동형사상을 \(K_\frakp/k_\frakp\)와 \((\Gg)_{k_\frakp}\) 사이의 \((\Gg)_{k_\frakp}\)-동형사상으로 확대할 수 있다. 이 동형사상은 \(\frakO_\frakp/\frako_\frakp\)에 대한 \([\Gg]_{\frako_\frakp}\)-동형사상을 유도한다. 여기서 \(K_\frakp/k_\frakp\)는 \(k_\frakp\) 위의 초복소계로 보아야 한다. 이는 \(K/k\)에서 계수영역을 확대하여 생기며, 일반적으로 영인자를 가진 계, 더 정확히는 서로 동형인 체들의 합인 반단순계가 된다.

\paragraph{정리 5.}
\emph{\(K/k\)의 차수 \(n\)을 나누지 않는 모든 소 \(\frakp\)에서 \(\frakO_\frakp/\frako_\frakp\)는 정규기저를 갖는다.}\footnote{아벨체의 경우에는 주 3)과 5)에 따라 국소화환으로도 소를 정의할 수 있다.}

정리 1에 따라 \([\Gg]_{\frako_\frakp}\)는 그곳에서 극대 오더이다. 따라서 \((\Gg)_{k_\frakp}\) 안에서 랭크 \(n\)인 \([\Gg]_{\frako_\frakp}\)-가군들은 정수 또는 분수 아이디얼이고, 정리 2에 따라 주아이디얼이다. 그러므로 \(\frakO_\frakp/\frako_\frakp\)에 대응하는 아이디얼이
\[
W[\Gg]_{\frako_\frakp}
\]
라면, 이는 \(\frako_\frakp\)-기저 \(WS_i\)를 갖는다. \(w\)가 \(\frakO_\frakp/\frako_\frakp\)에서 \(W\)에 대응하는 원소라 할 때, 연산자 동형사상은 \(w^{S_i}\)들이 \(\frakO_\frakp/\frako_\frakp\)의 \(\frako_\frakp\)-기저임을 말한다. 따라서 \(n\)개의 켤레원소 \(w^{S_i}\)가 \(\frakO_\frakp/\frako_\frakp\)의 정규기저를 이룬다.

\(\frakp\)가 \(n\)을 나누면 정규기저가 존재할 수 없으므로 \(\frakO_\frakp/\frako_\frakp\)에 대응하는 가군도 주가군일 수 없다.

\paragraph{정리 3에 대한 추가 설명.}
정리 3에서 증명한 \((\Gg)_k\)와 \(K/k\)의 연산자 동형사상\footnote{그 증명은 \(K\)가 \(k\) 위의 제1종 갈루아체이고 \(k\)가 무한히 많은 원소를 가진다는 것만을 가정한다. 이 마지막 제한은 불필요하다. M. Deuring은 유한체에도 유효한 정리 3의 증명을 고안하였고, 그 증명에서 거꾸로 체의 정규기저 존재가 따라온다. 동시에 유한체와 무한체에 한결같이 적용되는 원시원소 존재의 증명도 얻어진다.}에서 클라인의 형식문제에 관한 Speiser의 결과들\footnote{앞의 논문. 갈루아 가군의 개념은 그곳에서 추상화한 것이다. \(k\)의 표수가 \(n\)을 나누면 갈루아 가군에 따른 합성열과 그 기약인자들을 다루게 된다.}이 따라오며, 다음과 같이 정식화할 수 있다.

\emph{\(K/k\)의 갈루아군 \(\Gg\)의 \(k\) 위 기약표현은 모두 \(K/k\) 안의 한 기약 유리 갈루아 가군이 생성한다. 절대기약표현은 모두 \(K_Z/Z\) 안의 한 갈루아 가군이 생성한다.}

여기서 \(Z\)는 대응하는 군환 성분의 분해체로서 \(k\)를 포함한다. \(K_Z/Z\)는 \(Z\) 위의 초복소계로 보며 \(K/k\)의 계수영역 확대에서 생긴다. 특히 \(Z\)를 \(k\) 위에서 \(K\)와 선형적으로 서로소가 되게 택할 수 있으면 \(K_Z\)는 체로 남고 \(K_Z/Z\)는 부가확대가 된다. 이것이 Speiser가 다룬 경우이다.

명제 자체는 연산자 동형사상에 따라 \((\Gg)_k\) 또는 \((\Gg)_Z\)에 대한 대응 명제에서 따라온다. 그곳에서는 \((\Gg)_k\) 또는 \((\Gg)_Z\)에 관한 기약 아이디얼들이 표현을 생성한다.

\(v_1,\ldots,v_t\)가 그러한 갈루아 가군의 기저이고 \(S\mapsto\bar S\)가 대응하는 표현이면
\[
\begin{pmatrix}\vdots\\ v_i^S\\ \vdots\end{pmatrix}
=\bar S\begin{pmatrix}\vdots\\ v_i\\ \vdots\end{pmatrix}.
\]
이것이 Klein과 Speiser의 정식화이다.

\subsection*{§3. 고차 분기가 없는 체에서 군행렬식으로서의 판별식}

판별식을 분해하려면 각 분기소에서의 분해를 살펴보는 것으로 충분함이 알려져 있다. 따라서 고차 분기가 없는 체에서는 정규기저를 갖는 소들만을 다루면 된다.

\paragraph{정리 6.}
\emph{고차 분기가 없는 갈루아체 \(K/k\)에서는 각 분기소에서 판별식이 갈루아군의 군행렬식의 제곱과 같아지며, 이때 부정원들은 정규기저로 바뀐다. 절대기약표현들에 따라 군행렬식은 일반화된 근수들로 분해되고, 판별식은 지표들을 첨가한 기초체의 아이디얼들로 분해된다. 이때 복소켤레 지표들에는 같은 아이디얼들이 대응한다.}

\(w^S\)를 정규기저로 택하면 판별식 \(\Delta\)는
\[
D=\bigl|w^{ST^{-1}}\bigr|,\qquad S,T\in\Gg
\]
의 제곱과 같다. 그런데 \(D\)는 부정원 \(x_S\)들을 \(w^S\)로 바꾼 군행렬식이다.

근수들로의 분해는 앞에서 인용한 Speiser의 추론에서 얻어진다. 절대기약표현들에 따라 군행렬식과 군행렬을 각각
\[
D=D_1^{f_1}\cdots D_t^{f_t},
\qquad
M=f_1M_1+\cdots+f_tM_t
\]
로 분해하자. \(S\mapsto\bar S\)가 \(M_\lambda\)에 대응하는 표현이고 \(\bar S\)들이 확대된 기초체에 속하면
\[
M_\lambda=\sum_S w^S\bar S,
\qquad
M_\lambda^{T^{-1}}=\sum_S w^{ST^{-1}}\bar S
=\sum_R w^R\bar R\bar T=M_\lambda\bar T
\]
이고, 행렬식으로 넘어가면
\[
D_\lambda^{T^{-1}}=D_\lambda|\bar T|
=D_\lambda\varepsilon_T
\]
를 얻는다. 이로써 \(D_\lambda\)들이 일반화된 근수임을 알 수 있다. \(\varepsilon_T\)는 \(\bar T\)의 행렬식이므로 단위근이다. 이는 \(\Gg\)를 교환자부분군으로 나눈 몫군의 기약표현이다. 순환군의 경우 \(D_\lambda\)들은 단순히 라그랑주 근수
\[
\sum w^S\chi_\lambda(S)
\]
가 된다.

일반적으로 \(D_\lambda\)는 \(K_\frakp/k_\frakp\)에서 계수영역 \(k_\frakp\)에 대응하는 지표를 첨가하여 얻는 초복소계에 속하며, 더 정확히는 이 계의 \emph{정수적} 원소이다. 부정원 \(x_S\)들로 만든 행렬식의 계수들이 지표체의 대수적 정수이고 \(w^S\)들이 \(K_\frakp\)의 대수적 정수이기 때문이다.

군행렬식의 분해에서 판별식의 분해로 넘어가기 위해 매번 한 표현을 그 수반표현과 묶는다. \(\lambda,\bar\lambda\)가 서로 수반인 표현이면 동시에
\[
D_\lambda^{T^{-1}}=D_\lambda\varepsilon_T,
\qquad
D_{\bar\lambda}^{T^{-1}}=D_{\bar\lambda}\varepsilon_T^{-1}
\]
이다. 또 부정원 \(x_S\)들에 대하여 만든 \(\lambda,\bar\lambda\)의 행렬식들은 서로 복소켤레이고, 일반적으로 켤레 지표들에는 부정원 \(x_S\)들에 관한 켤레 행렬식들이 대응한다.

\(\Delta_\lambda=D_\lambda D_{\bar\lambda}\)로 두면 \(\lambda\)번째 지표체의 실부분체만 첨가하면 되고
\[
\Delta_\lambda^T=\Delta_\lambda
\qquad(\Gg\text{의 모든 }T)
\]
가 성립한다. 따라서\footnote{여기서는 한 체의 계수영역을 초복소적으로 확대하므로 통상적인 갈루아 이론의 추론을 수정해야 한다. \(P\)를 생각하는 \(k_\frakp\)의 확대체라 하고
\[
(K_\frakp)_P=(K_\frakp)_Pe^{S_1}+\cdots+(K_\frakp)_Pe^{S_r}
\]
를 체들로의 직합분해라 하자. 그러면 \((K_\frakp)_Pe^{S_i}/Pe^{S_i}\)는 갈루아가 되고, 그 군은 \(\frakp\)의 한 소인자의 분해군의 부분군이다. 또한 \(e^{S_i}\)들은 잉여류들에 관하여 켤레이다. \(\Delta_\lambda^T=\Delta_\lambda\)에서 먼저 \(\Delta_\lambda\)의 \(i\)번째 성분이 \(Pe^{S_i}\)에 속함이 따라오므로, 이는 \(P\)의 \(\gamma_i\)를 써서 \(\gamma_i e^{S_i}\)가 된다. 이어 \(e^{S_i}\)들의 켤레성에서 모든 \(\gamma_i\)가 같음이 따라온다. 따라서 실제로 \(\Delta_\lambda=\gamma\sum e^{S_i}=\gamma\)는 \(P\)에 속한다.} \(\Delta_\lambda\)는 \(\lambda\)번째 지표체의 실부분체를 첨가한 기초체에 속한다.

\emph{따라서 판별식 분해}
\[
\Delta=\Delta_1^{f_1}\cdots\Delta_t^{f_t}
\]
\emph{는 지표체들의 실부분체를 첨가한 기초체에서의 분해이며, 복소켤레 지표들에는 같은 인자들이 대응한다. 나머지 켤레 지표들에 대해서는 더 이상의 결론을 곧바로 내릴 수 없다[2a) 참조].}

이로써 정리 6의 모든 부분이 증명되었다. \(\Delta_\lambda\)들이 생성하는 아이디얼들이 실제로 이미 기초체 \(k\)에 속하고 켤레 지표들에 대하여 같아짐을 보일 수 있다면, 합성지표들에 대응하는 \(\Delta_\lambda\)들의 거듭제곱곱을 대응시키는 것만으로 이 아이디얼들은 Artin 도체들과 일치한다. 실제로 함수방정식이 성립하고, 정규기저에서 직접 따라오는 사실, 곧 부분군의 자명표현에는 대응하는 부분체의 판별식이 대응한다는 사실도 성립한다(앞서 인용한 Speiser 참조). 각 소에서 일치하면 전체 아이디얼도 일치한다. 유리수체 위에서 기약인 지표들로 제한하면 위 사실들만으로 일치가 따라온다. 절대기약지표에 대해서는 가장 간단한 경우에 그 일치를 직접 확인하겠다.

\paragraph{정리 7.}
\emph{\(k\)가 유리수체이고 \(K/k\)가 판별식과 서로소인 소수차수 \(l\)의 순환확대, 따라서 고차 분기가 없는 확대이면, 비자명표현들에 대하여 \(\Delta_\lambda\)들이 생성하는 아이디얼들은 서로 같고 \(K/k\)의 도체와 일치한다.}

이는 근수들의 알려진 분해에서 따라온다.\footnote{Hilbert, Zahlbericht, § 108. 정리 5에 따라 정규기저의 존재가 확정되어 있으므로 --- 주 7)에 따라 여기서는 소를 국소화환으로 정의해도 된다 --- 절대순환체가 원분체라는 Zahlbericht의 사실을 더는 끌어올 필요가 없다.} 실제로 \(K/k\)의 한 분기소 \(\frakp\)에서 \(l\)차 단위근체 \(k_\varepsilon\)와 \(K_\varepsilon\)에 대하여
\[
(p)=\frakp_1\cdots\frakp_{l-1},
\qquad
\frakp_i=\mathfrak P_i^{\,l}
\]
가 성립한다. 여기서 \(\frakp\)는 \(k_\varepsilon\)에, \(\mathfrak P\)는 \(K_\varepsilon\)에 속한다. \(k_\varepsilon\)는 \(K\)와 선형적으로 서로소이므로 \(K_\varepsilon\)는 체로 남고, 주 7)에 따라 여기서는 \(\frakp\)-진 확대가 필요 없다. 또한 근수들에 대하여
\[
(\Omega_\lambda)=
\mathfrak P_1^{r_1}\cdots\mathfrak P_{l-1}^{r_{l-1}},
\qquad
(\Omega_{\bar\lambda})=
\mathfrak P_1^{l-r_1}\cdots\mathfrak P_{l-1}^{l-r_{l-1}}
\]
가 성립하며 \(r_1,\ldots,r_{l-1}\)은 \(1,\ldots,l-1\)의 한 순열이다. 여기서는 근수 \(\Omega_\lambda\)가 군행렬식의 인자 \(D_\lambda\)와 동일하므로
\[
(\Delta_\lambda)=(D_\lambda D_{\bar\lambda})
=(\Omega_\lambda\Omega_{\bar\lambda})
=\mathfrak P_1^l\cdots\mathfrak P_{l-1}^l
=\frakp_1\cdots\frakp_{l-1}=(p)
\]
를 얻는다. 따라서 \(K/k\)의 도체와 일치한다.

\begin{center}
1931년 8월 24일 접수.
\end{center}
